\documentclass[11pt]{article}
\usepackage{graphicx}
\usepackage[
  backend=biber,
  style=ieee
]{biblatex}

\addbibresource{references.bib}

\title{AirJet Modeling}
\author{Nick Zhang}
\date{July 2026}

\begin{document}

\maketitle
\tableofcontents
\newpage

\section{Introduction}

\subsection{Context}

Nowadays electronic devices have better chips and more compact designs.
The power of chips keeps growing. However, most devices like laptops
either use passive cooling or a fan. These are the two classical
solutions, and each has clear trade-offs. A fan can handle higher power,
but it makes the device bigger, heavier, and noisier. Passive cooling
keeps the device light and quiet, but the cooling capacity is limited.
Because of this, even small mobile devices such as phones have started
using tiny fans. Conventional rotary-fan cooling has stayed the main
approach in most consumer electronics for a long time, while small
improvements have struggled to keep up with rising heat and shrinking
space~\cite{wang2013thermal}.

Even without increasing the total power, chips keep getting higher
transistor density. This makes the chip area smaller, which concentrates
the heat and makes cooling even harder~\cite{moore2014emerging}.

\subsection{AirJet Mini}

AirJet is a new way of active cooling that has no fans and no rotating
parts. It uses piezoelectrically actuated MEMS (microelectromechanical
systems) membranes vibrating at high frequency~\cite{hot_chips_2024}.
Each vibration cycle draws ambient air into internal cavities, pushes
it through arrays of microscopic orifices, and forms pulsating jets that
hit a heat-spreading surface. The heated air then leaves through a
lateral exhaust~\cite{frore_datasheet}.

The device measures 27.5 $\times$ 41.5 $\times$ 2.8~mm and weighs 11~g.
At the rated condition of 85~\textdegree{}C die temperature and
25~\textdegree{}C ambient temperature, it dissipates up to 5.25~W of
total heat while using at most 1~W of electrical power. Its maximum
backpressure is 1750~Pa and its in-system noise is 21~dBA measured at
50~cm~\cite{frore_datasheet}.

It uses very little space and makes less noise than a fan, so it fills
a niche between ordinary passive cooling and traditional fan-based
active cooling. We know what AirJet can do from its datasheet, but we
do not know how it is built inside. This paper tries to build a
complete model of the device using only public information.

\subsection{Prior Work}

AirJet has no rotating parts. Instead it uses high-frequency vibration
to create microjets. This mechanism, called a synthetic jet, has been
studied by researchers for many years. Past studies have looked at how
the distance from the orifice to the surface, written as the ratio H/D,
changes the cooling performance~\cite{krishan2019synthetic}.

Before AirJet, people also tried piezoelectric fans---a thin plate
driven by a piezoelectric element that swings back and forth to move
air~\cite{yoo2000piezoelectric}. These devices showed that piezoelectric
actuators can be used for cooling electronics, but they are open
structures with no enclosed cavity, no orifice plate, and no multi-cell
design.

More recently, some researchers built and tested a single-cavity
synthetic jet with multiple orifices for electronics cooling~\cite{singh2022experimental}.
This is close to a single AirJet cell. It shows that one cell can be
simulated and its cooling can be measured. However, no published study
has modeled a complete AirJet Mini with all cells, shared intake,
internal flow paths, and exhaust together.

Most of what we know about AirJet comes from the company itself: the
product datasheet, patents, and a Hot Chips 2024 tutorial. Independent
academic research on AirJet from outside Frore Systems is very rare.

\section{Methods}

\subsection{Evidence Sources and Classification}

To start the modeling, we first collect the available data and classify
them by how much we trust each piece.

The first and most trusted category comes directly from Frore Systems.
According to the manufacturer datasheet~\cite{frore_datasheet}, the
AirJet Mini Gen1 measures 27.5 $\times$ 41.5 $\times$ 2.8~mm and weighs
11~g. It uses up to 1~W of electrical power and dissipates 5.25~W of
total heat at 85~\textdegree{}C die temperature and 25~\textdegree{}C
ambient, which means 4.25~W of net heat removal from the chip. It can
work against a backpressure up to 1750~Pa, and its noise is 21~dBA at
50~cm.

The second category is Frore's patents~\cite{us12137540b2,us11978690b2}.
These documents give structural details as ranges instead of exact
numbers: for example, membrane diameters of 6--8~mm, membrane
displacements of 10--60~$\mu$m, upper cavity heights of 200--300~$\mu$m,
orifice widths of 200--300~$\mu$m, orifice spacing of 400--600~$\mu$m,
an open-area ratio of 8--12~\%, and an impingement gap of 200--300~$\mu$m.
Because the patents come directly from Frore Systems, these ranges are
trusted as reasonable bounds. In the model they are used as variables
that can be adjusted within physically sensible limits.

The third category is images from the datasheet and product page. The
cross-sectional diagram shows the general direction of intake and
exhaust, the rough internal zones (upper chamber, membrane, orifice
plate, impingement channel, and exhaust), and how the cells are likely
arranged inside the package. These images are manufacturer drawings, not
engineering blueprints. They are used only to guide layout choices; no
measurements are taken directly from them.

The fourth category is what we do not know. Several details are missing
from all public sources. The exact number of layers in the membrane, its
material, and its thickness are unknown. The piezoelectric material is
not named. The real driving frequency is not published---patent examples
mention 20--25~kHz for an 8~mm membrane, but this is just an example
range, not a confirmed product number. The real shape and size of
internal flow passages are unknown, and the heat-spreader material is
not identified. These unknowns are listed here because admitting what
we do not know is an important part of the approach.

\subsection{Geometry Reconstruction}

The external shape is fixed by the datasheet: 27.5 $\times$ 41.5
$\times$ 2.8~mm. For the internal structure, we only have a rough idea
from the cross-section diagram. From top to bottom, the device likely
contains: intake openings on the top surface, an upper chamber, a
vibrating membrane, a lower chamber, an orifice plate with arrays of
small holes, an impingement channel where jets hit the heat-spreader,
and a lateral exhaust~\cite{frore_datasheet,us12137540b2}.

The number of cells and their arrangement inside the package are not
publicly known. Several candidate layouts---with different cell counts,
cell sizes, and spacing---are considered in this study. All of them fit
within the fixed outer dimensions and follow the internal zoning shown
in patent drawings.

The total thickness of 2.8~mm must fit all internal layers. A thickness
budget splits this space across the upper chamber, membrane, lower
chamber, orifice plate, impingement channel, and heat-spreader. Ranges
for each layer thickness are taken from patent examples. Because these
exact numbers are not confirmed for the Mini Gen1, the budget is treated
as a set of limits, not a final answer.

\subsection{Multiscale Modeling Framework}

The AirJet Mini has multiple cells working at the same time. Simulating
every cell with full CFD---including moving membranes, compressible
flow, and heat transfer---would be too expensive for a first try.
Instead, we use a two-level approach.

We have no direct measurements for a single cell, but earlier
experimental and numerical studies on single-cavity synthetic jets~\cite{singh2022experimental,krishan2019synthetic}
have shown that this kind of cell-level simulation is possible and can
give reasonable flow and heat transfer results.

At the cell level, we build one high-detail model. It covers the full
cycle: the membrane moving up and down in a harmonic pattern, pressure
changes in the chambers, air speeding up through the orifice array, and
jet impingement on the heat-spreader surface. The output of this model
is the flow rate as a function of the pressure drop across the cell.

At the device level, each cell is treated as its own unit with the same
flow--pressure behavior obtained from the cell-level model. The
device-level model then calculates how the incoming flow splits among
cells, matches each cell to the shared exhaust, and gives the total
heat removed from the heat-spreader surface.

This approach makes the full-device simulation practical. The main
limitation is that interaction between cells---like acoustic coupling
between neighboring chambers---is not yet included in the cell-level
model.

\subsection{Boundary Conditions and Assumptions}

The main case uses the rated condition from the datasheet: the chip at
85~\textdegree{}C and the surrounding air at 25~\textdegree{}C~\cite{frore_datasheet}.
We test two exhaust conditions: open discharge at 0~Pa, and the maximum
backpressure of 1750~Pa. These are the easiest and hardest conditions,
and we expect the real operating point to fall somewhere between them.

The membrane moves in a sinusoidal, harmonic pattern, matching the
piezoelectric drive described in the patents~\cite{us12137540b2}. We
do not lock in the exact frequency and displacement yet. These are
taken from patent ranges and will be decided later through calibration.

For the first simulation, air is treated as incompressible with
constant properties at 25~\textdegree{}C. This keeps things simple and
lets us get a working flow field first. Compressibility can be added
afterwards.

The heat-conducting bottom receives 4.25~W, which is the heat from the
chip. The AirJet itself uses up to 1~W of electrical power, and this
becomes heat too. We put it as a separate heat source in the actuator
area. Together they add up to the 5.25~W total dissipation from the
datasheet. The outside walls of the package are treated as adiabatic
for now.

A few effects are left out at this stage: radiation, acoustic coupling
between neighboring cells, and exact material properties. We note them
here and they can be added later.

\printbibliography

\end{document}
